\scanpage{4}
\mspage{2}

% opening
$l-8$\qquad pourquoi \uline{donc}\qquad ?

% division
\hrule

% p 3
p 3\quad $l-8,-9$\qquad $B$ ?

% division
\hrule

% p 4, 1.3.1
p 4.\quad 1.3.1.\qquad As tu vérifi\deleted{\unsure{e}}é si l'accouplement d'A. M. coïncide avec le tien ?

% p 4, l 4
$l\,4$.\qquad \uline{formellement} non ramifié ?

% p 4, l 7
$l\,7$.\qquad \uline{séparablement} clos ?

% p 4, l 14–20
$l\,14-20$.\quad Est-ce utile à dire, en l'absence de référence à ce que ce peut vouloir dire ?

% p 4, l −5
$l-5$\qquad \unsure{Qu'il|qu'il} soit \qquad et \qquad signifierait \qquad plutôt ?

% division
\hrule

% pg 5
pg 5\quad $l-2$\qquad français :\quad les classes $\unsure{w|\omega}$\qquad \uline{ou}\qquad relatifs.

% division
\hrule

% pg 6
pg 6\quad $l\,4$\qquad français. dire\qquad un Module inversible ample\qquad ?

% pg 6, below
\deleted{\illegible{4}}

% division
\hrule\smallskip\hrule\smallskip\hrule

% section
\tikz[baseline=(c.base)]{\node[draw,double,ellipse,inner sep=1pt] (c) {2};}

% pg 7
pg 7\qquad français \uline{$l-9$}\quad \tikz[baseline=(c.base)]{\node[draw,ellipse,inner sep=1pt] (c) {$\unsure{c|\subset}$};}$L_0$\quad donc \tikz[baseline=(c.base)]{\node[draw,ellipse,inner sep=1pt] (c) {c'};} est

% pg 7, second line
\deleted{$l- \unsure{8|7|1}$ \unsure{si} \unsure{que}}\quad $l-7$\quad \tikz[baseline=(c.base)]{\node[draw,ellipse,inner sep=1pt] (c) {que};}\quad (le $2^{\underline{\mathrm o}}$)

% division
\hrule

% pg 8
pg 8\qquad français\quad $l\,4$\qquad par exemple \inserted{\tikz[baseline=(c.base)]{\node[draw,ellipse,inner sep=1pt] (c) {pour};} ?}\unsure{k|si} $\ldots$

% pg 8, next line
$l-7$\quad \uline{le} terme

% division
\hrule

% pg 14
pg 14.\qquad L'argument est un peu immoral ; si $X$ et $Y$ sont extension d'une VA par un tore et si $\varphi:X\longrightarrow Y$ est donné, alors $\varphi$ est strictement compatible à la filtration par le poids et $\varphi$ est une isogénie si et seulement si $T_\ell(\varphi)$ en est une.

% division
\hrule

% pg 20
\tikz[baseline=(c.base)]{\node[draw,ellipse,inner sep=1pt] (c) {pg \deleted{\unsure{15|18}} 20.};}\quad Rmq 2.6 a)\qquad $S$ est-il un trait dans cette remarque ; sait-on assez sur le passage à la limite et Néron ?
% archive sheet number: 3
